% ============================================================
%  C: Como Programar -- 6ª Edição | Deitel & Deitel
%  Capítulo 3 -- Desenvolvimento Estruturado de Programas em C
%  EXERCÍCIOS DE AUTORREVISÃO (3.1 -- 3.10)
%  Abordagem incremental: Caps. 1 + 2 + 3
% ============================================================
\documentclass[12pt,a4paper]{article}
\usepackage[utf8]{inputenc}
\usepackage[T1]{fontenc}
\usepackage[brazil]{babel}
\usepackage{geometry}
\geometry{top=2.5cm,bottom=2.5cm,left=3cm,right=3cm}
\usepackage{parskip}\setlength{\parskip}{6pt}\setlength{\parindent}{0pt}
\usepackage{xcolor,tcolorbox,enumitem,listings,fancyhdr,titlesec,hyperref,booktabs,array}
\tcbuselibrary{skins,breakable}

\definecolor{deitelblue}{RGB}{0,70,127}
\definecolor{deitelgray}{RGB}{240,242,245}
\definecolor{deitelgreen}{RGB}{0,110,60}
\definecolor{deitellightblue}{RGB}{220,235,250}
\definecolor{deitelred}{RGB}{160,20,20}
\definecolor{deitellorange}{RGB}{200,90,0}

\newtcolorbox{boaprática}[1][]{colback=deitellightblue,colframe=deitelblue,
  fonttitle=\bfseries\small,title={Boa prática de programação},breakable,#1}
\newtcolorbox{erroco}[1][]{colback=red!5,colframe=deitelred,
  fonttitle=\bfseries\small,title={Erro comum de programação},breakable,#1}
\newtcolorbox{prev}[1][]{colback=yellow!8,colframe=deitellorange,
  fonttitle=\bfseries\small,title={Dica de prevenção de erro},breakable,#1}
\newtcolorbox{respostabox}[1][]{colback=deitelgray,colframe=deitelblue!60,
  fonttitle=\bfseries\small\color{deitelblue},title={Resposta detalhada},breakable,#1}
\newtcolorbox{enunciadoBox}[1][]{colback=white,colframe=deitelblue,
  fonttitle=\bfseries,title={#1},breakable}
\newtcolorbox{saida}[1][]{colback=black!88,colframe=deitelblue,
  fontupper=\ttfamily\small\color{green!70},title={\small\color{white}Saída do programa},breakable,#1}

\setlist[enumerate,1]{label=\textbf{\alph*)},leftmargin=2em}
\setlist[itemize]{leftmargin=2em}

\lstset{
  basicstyle=\ttfamily\small,backgroundcolor=\color{deitelgray},
  frame=single,framerule=0.4pt,rulecolor=\color{deitelblue!40},
  breaklines=true,showstringspaces=false,
  keywordstyle=\color{deitelblue}\bfseries,
  commentstyle=\color{deitelgreen}\itshape,
  stringstyle=\color{deitelred},
  language=C,numbers=left,numberstyle=\tiny\color{gray},
  xleftmargin=18pt,xrightmargin=5pt,stepnumber=1,
}

\pagestyle{fancy}\fancyhf{}
\fancyhead[L]{\small\color{deitelblue}\textbf{C: Como Programar -- 6ª ed. | Deitel \& Deitel}}
\fancyhead[R]{\small\color{deitelblue}Cap. 3 -- Autorrevisão}
\fancyfoot[C]{\small\thepage}
\renewcommand{\headrulewidth}{0.4pt}

\titleformat{\section}[block]{\large\bfseries\color{deitelblue}}{\thesection.}{0.5em}{}[\titlerule]
\titleformat{\subsection}[block]{\normalsize\bfseries\color{deitelblue!80}}{}{}{}

\hypersetup{colorlinks=true,linkcolor=deitelblue}

\begin{document}

\begin{titlepage}
  \centering\vspace*{2cm}
  {\color{deitelblue}\rule{\linewidth}{2pt}}\\[0.4cm]
  {\huge\bfseries\color{deitelblue}C: Como Programar}\\[0.2cm]
  {\large\color{deitelblue}6ª Edição $\cdot$ Paul Deitel \& Harvey Deitel}\\[0.2cm]
  {\color{deitelblue}\rule{\linewidth}{2pt}}\\[1cm]
  {\LARGE\bfseries Capítulo 3}\\[0.3cm]
  {\Large\color{deitelblue!80}Desenvolvimento Estruturado de Programas em C}\\[0.8cm]
  \begin{tcolorbox}[colback=deitellightblue,colframe=deitelblue,width=0.85\linewidth]
    \centering{\large\bfseries Exercícios de Autorrevisão (3.1--3.10)}\\[0.2cm]
    {\normalsize Resolvidos passo a passo, com abordagem incremental\\
    Capítulos 1 + 2 + 3}
  \end{tcolorbox}
  \vfill{\small Abordagem incremental Deitel}
\end{titlepage}

\section*{Construindo sobre os Capítulos Anteriores}

No Capítulo~3, ampliamos consideravelmente nosso arsenal de ferramentas em C.
Construindo sobre os fundamentos dos Capítulos~1 e~2, incorporamos agora:

\begin{itemize}[noitemsep]
  \item \textbf{Algoritmos e pseudocódigo} -- a etapa de planejamento que precede a codificação
  \item \textbf{Refinamentos top-down} -- a metodologia para projetar programas estruturados
  \item \textbf{\texttt{if...else}} -- estrutura de seleção dupla
  \item \textbf{\texttt{while}} -- nossa primeira estrutura de repetição
  \item \textbf{Repetição controlada por contador} (definida) e \textbf{por sentinela} (indefinida)
  \item \textbf{Tipo \texttt{float}} e cálculos com ponto flutuante
  \item \textbf{Operadores de atribuição compostos} (\texttt{+=}, \texttt{-=}, \texttt{*=}, \texttt{/=}, \texttt{\%=})
  \item \textbf{Operadores de incremento/decremento} (\texttt{++} e \texttt{--}) -- prefixo e pós-fixo
\end{itemize}

\begin{boaprática}
  Os conceitos do Capítulo~3 são o coração da programação estruturada.
  Domine pseudocódigo, refinamentos top-down e \texttt{while} -- esses três
  elementos formam a base sobre a qual todo programa não-trivial será construído.
\end{boaprática}

\tableofcontents\newpage

% ============================================================
\section{Exercício 3.1 -- Preenchimento de Lacunas}
% ============================================================

\begin{enunciadoBox}[Enunciado 3.1]
Preencha os espaços nas oito sentenças sobre algoritmos, controle de programa,
estruturas de seleção, instruções compostas e estruturas de repetição.
\end{enunciadoBox}

\subsection*{Item (a) -- Procedimento para resolver problema em termos de ações ordenadas}

\begin{respostabox}
\textbf{Resposta:} \textcolor{deitelblue}{\textbf{algoritmo}}

\medskip\textbf{Explicação:}

Da Seção~3.2: um \textbf{algoritmo} é o procedimento utilizado para resolver um problema
em termos (1) das ações a serem executadas e (2) da ordem em que essas ações devem
ser executadas. A ordem é fundamental -- como ilustramos no exemplo do executivo
``preparando-se para o trabalho'': se ele veste-se antes de tomar banho, o resultado
será desastroso! Em programação, especificar a ordem correta dos passos é a essência
de qualquer algoritmo.
\end{respostabox}

\subsection*{Item (b) -- Especificação da ordem de execução das instruções}

\begin{respostabox}
\textbf{Resposta:} \textcolor{deitelblue}{\textbf{controle do programa}}

\medskip\textbf{Explicação:}

\textbf{Controle do programa} é o termo técnico para a especificação da ordem em
que as instruções de um programa devem ser executadas. O Capítulo~3 dedica-se quase
integralmente a estudar os mecanismos de controle do programa em C: as estruturas
de seleção (\texttt{if}, \texttt{if...else}) e as estruturas de repetição (\texttt{while}).
\end{respostabox}

\subsection*{Item (c) -- Os três tipos de estruturas de controle}

\begin{respostabox}
\textbf{Resposta:} \textcolor{deitelblue}{\textbf{sequência, seleção e repetição}}

\medskip\textbf{Explicação:}

A pesquisa de Bohm e Jacopini (1966) demonstrou um resultado notável: \textit{qualquer
programa pode ser escrito utilizando apenas três tipos de estruturas de controle}:

\begin{enumerate}
  \item \textbf{Sequência:} As instruções são executadas uma após a outra, na ordem
  em que aparecem. Esta é a estrutura ``default'' -- está embutida na linguagem C.
  \item \textbf{Seleção:} O programa escolhe entre cursos alternativos de ação.
  Em C: \texttt{if} (seleção única), \texttt{if...else} (seleção dupla),
  \texttt{switch} (seleção múltipla, Cap.~4).
  \item \textbf{Repetição:} Um conjunto de instruções é executado várias vezes
  enquanto uma condição permanece verdadeira. Em C: \texttt{while},
  \texttt{do...while} e \texttt{for} (estes dois últimos no Cap.~4).
\end{enumerate}

Essa descoberta é a base teórica da \textbf{programação estruturada}, que tornou
quase sinônima a expressão ``eliminação do \texttt{goto}''.
\end{respostabox}

\subsection*{Item (d) -- Estrutura de seleção que executa uma ação ou outra}

\begin{respostabox}
\textbf{Resposta:} \textcolor{deitelblue}{\textbf{\texttt{if...else}}}

\medskip\textbf{Explicação:}

A estrutura \texttt{if...else} (Seção~3.6) é chamada de \textit{seleção dupla}:
ela executa uma ação se a condição for verdadeira e outra ação diferente se a
condição for falsa. Sintaxe básica:

\begin{lstlisting}
if ( nota >= 60 ) {
    printf( "Aprovado\n" );
} /* fim do if */
else {
    printf( "Reprovado\n" );
} /* fim do else */
\end{lstlisting}

Diferentemente do \texttt{if} simples (Cap.~2), que apenas executa ou não uma ação,
o \texttt{if...else} sempre executa \textit{exatamente um} dos dois caminhos.
\end{respostabox}

\subsection*{Item (e) -- Várias instruções agrupadas com chaves}

\begin{respostabox}
\textbf{Resposta:} \textcolor{deitelblue}{\textbf{instrução composta}} (ou \textbf{bloco})

\medskip\textbf{Explicação:}

Um conjunto de instruções entre chaves \texttt{\{} e \texttt{\}} forma uma
\textbf{instrução composta} (também chamada de \textbf{bloco}). Uma instrução composta
pode aparecer em qualquer lugar onde uma instrução única seria válida -- esse é
um conceito poderoso da linguagem C.

\begin{lstlisting}
if ( nota >= 60 ) {
    printf( "Aprovado\n" );        /* instrucao 1 do bloco */
    printf( "Parabens!\n" );       /* instrucao 2 do bloco */
    aprovados = aprovados + 1;     /* instrucao 3 do bloco */
} /* todas executadas se nota >= 60 */
\end{lstlisting}
\end{respostabox}

\subsection*{Item (f) -- Estrutura de repetição com condição verdadeira}

\begin{respostabox}
\textbf{Resposta:} \textcolor{deitelblue}{\textbf{\texttt{while}}}

\medskip\textbf{Explicação:}

A estrutura \texttt{while} (Seção~3.7) repete uma instrução ou bloco
\textit{enquanto} uma condição permanecer verdadeira. Quando a condição se torna
falsa, o laço termina e a execução continua na próxima instrução após o
\texttt{while}.

\begin{lstlisting}
contador = 1;
while ( contador <= 10 ) {       /* enquanto contador <= 10 */
    printf( "%d\n", contador );  /* exibe o valor          */
    contador = contador + 1;     /* incrementa o contador  */
} /* fim do while */
\end{lstlisting}

\begin{erroco}
  Esquecer de alterar a variável da condição dentro do corpo do \texttt{while}
  é o erro mais clássico associado a essa estrutura: cria-se um \textbf{loop
  infinito}. Sempre garanta que algum elemento do corpo do laço modifique a
  condição em direção a se tornar falsa.
\end{erroco}
\end{respostabox}

\subsection*{Item (g) -- Repetição por número específico de vezes}

\begin{respostabox}
\textbf{Resposta:} \textcolor{deitelblue}{\textbf{controlada por contador}} (ou \textbf{repetição definida})

\medskip\textbf{Explicação:}

A \textbf{repetição controlada por contador} (Seção~3.8) usa uma variável-contador
para especificar exatamente quantas vezes um conjunto de instruções deve ser executado.
É chamada de \textit{repetição definida} porque o número de repetições é
\textbf{conhecido antes} do laço começar.

Padrão típico:
\begin{enumerate}[noitemsep]
  \item Inicialize o contador (geralmente com 1)
  \item No \texttt{while}, teste se o contador atingiu o limite
  \item No corpo, execute a tarefa e incremente o contador
\end{enumerate}
\end{respostabox}

\subsection*{Item (h) -- Repetição quando o número de iterações é desconhecido}

\begin{respostabox}
\textbf{Resposta:} \textcolor{deitelblue}{\textbf{sentinela}} (também chamado de
\textit{sinalizador}, \textit{valor artificial} ou \textit{flag})

\medskip\textbf{Explicação:}

Quando não sabemos com antecedência quantas vezes o laço será executado, usamos um
\textbf{valor de sentinela} -- um valor especial que o usuário digita para indicar
o ``fim da entrada de dados''. A repetição controlada por sentinela é também chamada
de \textit{repetição indefinida}.

\textbf{Regra de ouro:} O valor de sentinela deve ser escolhido de modo que
\textit{nunca} possa ser confundido com um dado válido. Por exemplo, em um programa
que lê notas de teste (sempre $\geq 0$), \texttt{-1} é uma boa sentinela.
\end{respostabox}

\newpage

% ============================================================
\section{Exercício 3.2 -- Quatro formas de incrementar \texttt{x}}
% ============================================================

\begin{enunciadoBox}[Enunciado 3.2]
Escreva quatro instruções diferentes em C; cada uma deve somar 1 à variável inteira \texttt{x}.
\end{enunciadoBox}

\begin{respostabox}
C oferece múltiplas formas para a operação fundamental de incrementar uma variável.
Cada uma delas tem o \textit{mesmo efeito} (somar 1 a \texttt{x}) quando usada
isoladamente como uma instrução:

\begin{lstlisting}
x = x + 1;   /* forma classica usando atribuicao e adicao */
x += 1;      /* operador de atribuicao composto +=         */
++x;         /* operador de pre-incremento                 */
x++;         /* operador de pos-incremento                 */
\end{lstlisting}

\textbf{Quando usar cada forma?}
\begin{itemize}
  \item \texttt{x = x + 1} -- forma explícita, mais legível para iniciantes
  \item \texttt{x += 1} -- forma compacta, útil para incrementos por valor diferente de 1
  \item \texttt{++x} ou \texttt{x++} -- forma idiomática em C, especialmente em laços
\end{itemize}

\begin{boaprática}
  Embora as quatro formas sejam equivalentes \textit{em uma instrução isolada},
  o pré-incremento (\texttt{++x}) e o pós-incremento (\texttt{x++}) podem ter
  efeitos diferentes quando usados \textit{dentro de uma expressão maior} -- como
  veremos no Exercício~3.6.
\end{boaprática}
\end{respostabox}

\newpage

% ============================================================
\section{Exercício 3.3 -- Instruções únicas em C}
% ============================================================

\begin{respostabox}
\begin{enumerate}[label=\textbf{\alph*)}]

\item \textbf{Atribuir \texttt{x + y} a \texttt{z} e incrementar \texttt{x} após o cálculo:}
\begin{lstlisting}
z = x++ + y;
\end{lstlisting}
O pós-incremento \texttt{x++} usa o valor \textit{atual} de \texttt{x} na soma e
\textit{depois} incrementa \texttt{x} -- exatamente o comportamento solicitado.

\item \textbf{Multiplicar \texttt{produto} por 2 com \texttt{*=}:}
\begin{lstlisting}
produto *= 2;
\end{lstlisting}
Equivalente a \texttt{produto = produto * 2}.

\item \textbf{Multiplicar \texttt{produto} por 2 com \texttt{=} e \texttt{*}:}
\begin{lstlisting}
produto = produto * 2;
\end{lstlisting}
Forma explícita, sem o operador composto.

\item \textbf{Testar se \texttt{contador > 10} e exibir mensagem:}
\begin{lstlisting}
if ( contador > 10 ) {
    printf( "Contador e maior do que 10\n" );
} /* fim do if */
\end{lstlisting}

\item \textbf{Decrementar \texttt{x} em 1, depois subtrair de \texttt{total}:}
\begin{lstlisting}
total -= --x;
\end{lstlisting}
O pré-decremento \texttt{--x} primeiro decrementa \texttt{x} e \textit{depois} usa o
novo valor na subtração: \texttt{total = total - (x - 1)}, com \texttt{x} já alterado.

\item \textbf{Somar \texttt{x} a \texttt{total}, depois decrementar \texttt{x}:}
\begin{lstlisting}
total += x--;
\end{lstlisting}
O pós-decremento \texttt{x--} primeiro usa o valor atual de \texttt{x} na soma e
\textit{depois} decrementa \texttt{x}.

\item \textbf{Resto de \texttt{q / divisor} atribuído a \texttt{q} (duas formas):}
\begin{lstlisting}
q %= divisor;        /* forma compacta com %=  */
q = q % divisor;     /* forma explicita        */
\end{lstlisting}

\item \textbf{Imprimir \texttt{123.4567} com 2 dígitos de precisão:}
\begin{lstlisting}
printf( "%.2f", 123.4567 );
\end{lstlisting}
\textbf{Saída:} \texttt{123.46} -- o valor é \textit{arredondado} (não truncado)
para duas casas decimais. Como o terceiro dígito é 6 ($\geq 5$), arredonda-se
para cima: \texttt{.45 \(\to\) .46}.

\item \textbf{Imprimir \texttt{3.14159} com 3 dígitos à direita do ponto decimal:}
\begin{lstlisting}
printf( "%.3f\n", 3.14159 );
\end{lstlisting}
\textbf{Saída:} \texttt{3.142} -- arredondado (o quarto dígito, 5, força o arredondamento
para cima: \texttt{.141 \(\to\) .142}, na verdade o que acontece é que se olha o próximo
dígito e arredonda-se).

\end{enumerate}
\end{respostabox}

\newpage

% ============================================================
\section{Exercício 3.4 -- Instruções para soma}
% ============================================================

\begin{respostabox}
\begin{enumerate}[label=\textbf{\alph*)}]

\item \textbf{Declarar \texttt{soma} e \texttt{x} como \texttt{int}:}
\begin{lstlisting}
int soma, x;
\end{lstlisting}

\item \textbf{Inicializar \texttt{x} em 1:}
\begin{lstlisting}
x = 1;
\end{lstlisting}

\item \textbf{Inicializar \texttt{soma} em 0:}
\begin{lstlisting}
soma = 0;
\end{lstlisting}

\item \textbf{Somar \texttt{x} a \texttt{soma}:}
\begin{lstlisting}
soma += x;       /* forma compacta   */
soma = soma + x; /* forma equivalente */
\end{lstlisting}

\item \textbf{Imprimir mensagem com valor de \texttt{soma}:}
\begin{lstlisting}
printf( "A soma e: %d\n", soma );
\end{lstlisting}

\end{enumerate}

\begin{boaprática}
  Observe a regra fundamental do Cap.~3: \textbf{contadores e totais devem ser
  inicializados antes do uso}. Variáveis não inicializadas contêm ``lixo'' --
  o último valor que ocupava aquela posição de memória, totalmente imprevisível.
\end{boaprática}
\end{respostabox}

\newpage

% ============================================================
\section{Exercício 3.5 -- Soma dos inteiros de 1 a 10}
% ============================================================

\begin{enunciadoBox}[Enunciado 3.5]
Combine as instruções do Exercício 3.4 em um programa completo que calcule a soma
dos inteiros de 1 a 10, usando a estrutura \texttt{while}. O laço deve terminar
quando \texttt{x} atingir 11.
\end{enunciadoBox}

\begin{respostabox}
\begin{lstlisting}
/* Calcula a soma dos inteiros de 1 a 10 */
#include <stdio.h>

int main( void )
{
    int soma, x; /* declara variaveis soma e x */

    x = 1;       /* inicializa x */
    soma = 0;    /* inicializa soma */

    while ( x <= 10 ) {  /* laco enquanto x menor ou igual a 10 */
        soma += x;       /* soma x ao total acumulado            */
        ++x;             /* incrementa x                         */
    } /* fim do while */

    printf( "A soma e: %d\n", soma );
    return 0;
} /* fim da funcao main */
\end{lstlisting}

\begin{saida}
A soma e: 55
\end{saida}

\medskip
\textbf{Rastreamento (trace) das primeiras iterações:}

\begin{center}
\renewcommand{\arraystretch}{1.4}
\begin{tabular}{c c c c}
  \toprule
  \textbf{Iteração} & \textbf{x (início)} & \textbf{soma (após \texttt{soma += x})} & \textbf{x (após \texttt{++x})} \\ \midrule
  Inicialização & 1 & 0 & -- \\
  1 & 1 & 1 & 2 \\
  2 & 2 & 3 & 3 \\
  3 & 3 & 6 & 4 \\
  4 & 4 & 10 & 5 \\
  $\vdots$ & $\vdots$ & $\vdots$ & $\vdots$ \\
  10 & 10 & 55 & 11 \\
  Saída do laço & 11 & 55 & -- \\
  \bottomrule
\end{tabular}
\end{center}

A condição \texttt{11 <= 10} é falsa, e o laço termina. \texttt{soma} contém a
resposta correta: $1+2+\ldots+10 = 55$.
\end{respostabox}

\newpage

% ============================================================
\section{Exercício 3.6 -- Pós-incremento dentro de expressão}
% ============================================================

\begin{enunciadoBox}[Enunciado 3.6]
Determine os valores de \texttt{produto} e \texttt{x} após a execução de
\texttt{produto *= x++;}, supondo que ambas começam com valor 5.
\end{enunciadoBox}

\begin{respostabox}
\textbf{Resposta:} \textcolor{deitelblue}{\textbf{produto = 25, x = 6}}

\medskip\textbf{Explicação passo a passo:}

A expressão \texttt{produto *= x++;} é equivalente a
\texttt{produto = produto * x++;}. O comportamento do \textbf{pós-incremento} é
crítico aqui:

\begin{enumerate}
  \item \texttt{x++} primeiro \textbf{usa o valor atual de \texttt{x}} (que é 5)
  na expressão \texttt{produto * x}.
  \item Calcula-se \texttt{produto * 5 = 5 * 5 = 25}, que é atribuído a \texttt{produto}.
  \item \textbf{Após} o uso do valor, \texttt{x} é incrementado para 6.
\end{enumerate}

\textbf{Resultado final:} \texttt{produto = 25}, \texttt{x = 6}.

\medskip
\textbf{Comparação com o pré-incremento:}

Se a instrução fosse \texttt{produto *= ++x;}, o resultado seria diferente:
\begin{enumerate}[noitemsep]
  \item \texttt{++x} primeiro \textbf{incrementa \texttt{x} para 6}.
  \item Depois, calcula-se \texttt{produto * 6 = 5 * 6 = 30}.
  \item Resultado: \texttt{produto = 30}, \texttt{x = 6}.
\end{enumerate}

\begin{prev}
  Esta diferença sutil entre \texttt{x++} e \texttt{++x} \textit{dentro de uma
  expressão} é uma fonte clássica de bugs em C. Quando estiver em dúvida sobre
  o comportamento, separe a operação em duas instruções para tornar o código
  inequívoco.
\end{prev}
\end{respostabox}

\newpage

% ============================================================
\section{Exercício 3.7 -- Instruções para potência}
% ============================================================

\begin{respostabox}
\begin{enumerate}[label=\textbf{\alph*)}]
\item \texttt{scanf( "\%d", \&x );}
\item \texttt{scanf( "\%d", \&y );}
\item \texttt{i = 1;}
\item \texttt{potencia = 1;}
\item \texttt{potencia *= x;} \quad (ou \texttt{potencia = potencia * x;})
\item \texttt{i++;} \quad (ou \texttt{++i;}, ou \texttt{i = i + 1;}, ou \texttt{i += 1;})
\item \texttt{while ( i <= y )}
\item \texttt{printf( "\%d", potencia );}
\end{enumerate}
\end{respostabox}

% ============================================================
\section{Exercício 3.8 -- Programa: $x$ elevado a $y$}
% ============================================================

\begin{enunciadoBox}[Enunciado 3.8]
Use as instruções do Exercício 3.7 para calcular $x^y$ usando \texttt{while}.
\end{enunciadoBox}

\begin{respostabox}
\begin{lstlisting}
/* Eleva x a potencia y */
#include <stdio.h>

int main( void )
{
    int x, y, i, potencia;  /* declara variaveis */

    i = 1;          /* inicializa contador      */
    potencia = 1;   /* inicializa potencia em 1 */

    printf( "Digite a base x: " );
    scanf( "%d", &x );

    printf( "Digite o expoente y: " );
    scanf( "%d", &y );

    while ( i <= y ) {       /* enquanto i <= y                */
        potencia *= x;       /* multiplica potencia por x      */
        ++i;                 /* incrementa o contador          */
    } /* fim do while */

    printf( "%d elevado a %d e %d\n", x, y, potencia );

    return 0;
} /* fim da funcao main */
\end{lstlisting}

\begin{saida}
Digite a base x: 2\\
Digite o expoente y: 10\\
2 elevado a 10 e 1024
\end{saida}

\textbf{Como funciona:} multiplicamos \texttt{potencia} por \texttt{x} repetidamente,
\texttt{y} vezes. Após \texttt{y} iterações, \texttt{potencia} contém $x^y$.
Por exemplo, $2^{10}$: $1 \to 2 \to 4 \to 8 \to 16 \to 32 \to 64 \to 128 \to
256 \to 512 \to 1024$.
\end{respostabox}

\newpage

% ============================================================
\section{Exercício 3.9 -- Identificar e Corrigir Erros}
% ============================================================

\begin{respostabox}
\begin{enumerate}[label=\textbf{\alph*)}]

\item \textbf{Código com erro:}
\begin{lstlisting}
while ( c <= 5 ) {
    produto *= c;
    ++c;
\end{lstlisting}
\textbf{Erro:} Falta a chave de fechamento \texttt{\}} do \texttt{while}.\\
\textbf{Correção:}
\begin{lstlisting}
while ( c <= 5 ) {
    produto *= c;
    ++c;
} /* fim do while */
\end{lstlisting}

\item \textbf{Código com erro:}
\begin{lstlisting}
scanf( "%.4f", &valor );
\end{lstlisting}
\textbf{Erro:} A precisão (\texttt{.4}) não é válida em uma especificação de conversão
do \texttt{scanf}. A precisão é usada apenas em \texttt{printf} para controlar a
formatação da saída.\\
\textbf{Correção:}
\begin{lstlisting}
scanf( "%f", &valor );
\end{lstlisting}

\item \textbf{Código com erro:}
\begin{lstlisting}
if ( sexo == 1 )
    printf( "Mulher\n" );
se nao
    printf( "Homem\n" );
\end{lstlisting}
\textbf{Erro:} A palavra \texttt{se nao} (em português) não existe em C. A palavra-chave
correta é \texttt{else}.\\
\textbf{Correção:}
\begin{lstlisting}
if ( sexo == 1 ) {
    printf( "Mulher\n" );
}
else {
    printf( "Homem\n" );
}
\end{lstlisting}

\end{enumerate}
\end{respostabox}

% ============================================================
\section{Exercício 3.10 -- Loop Infinito}
% ============================================================

\begin{enunciadoBox}[Enunciado 3.10]
Identifique o erro no \texttt{while} a seguir (\texttt{z} começa em 100), que
deveria calcular a soma dos inteiros de 100 a 1:
\begin{lstlisting}
while ( z >= 0 )
    soma += z;
\end{lstlisting}
\end{enunciadoBox}

\begin{respostabox}
\textbf{Erro:} O valor de \texttt{z} \textbf{nunca é alterado} dentro do corpo do
\texttt{while}. Como \texttt{z} permanece sempre 100, a condição \texttt{z >= 0}
nunca se torna falsa, e o programa entra em um \textbf{loop infinito}.

\medskip\textbf{Correção:} é necessário decrementar \texttt{z} dentro do laço,
para que ele eventualmente chegue a $-1$ e a condição se torne falsa:

\begin{lstlisting}
while ( z >= 0 ) {
    soma += z;
    --z;          /* decrementa z para evitar loop infinito */
} /* fim do while */
\end{lstlisting}

Com essa correção, \texttt{z} assume os valores 100, 99, 98, \ldots, 1, 0, e
\texttt{soma} acumula \(100+99+\ldots+1+0 = 5050\). Quando \texttt{z} chega a $-1$,
a condição falha e o laço termina.

\begin{erroco}
  \textbf{Loop infinito} é um dos erros mais comuns em programação. A regra é
  simples: o corpo do laço \textit{deve} alterar, eventualmente, alguma variável
  da condição em direção a torná-la falsa. Se isso não acontecer, o programa
  ficará rodando indefinidamente, exigindo interrupção forçada (Ctrl+C).
\end{erroco}
\end{respostabox}

\newpage

% ============================================================
\section*{Gabarito Rápido -- Capítulo 3: Autorrevisão}

\begin{tcolorbox}[colback=deitellightblue,colframe=deitelblue,title={\bfseries\large Respostas Resumidas}]

\textbf{3.1}
\begin{itemize}[noitemsep]
  \item[a)] algoritmo \quad b) controle do programa
  \item[c)] sequência, seleção, repetição \quad d) \texttt{if...else}
  \item[e)] instrução composta (bloco) \quad f) \texttt{while}
  \item[g)] controlada por contador \quad h) sentinela
\end{itemize}

\medskip\textbf{3.2:} \texttt{x = x + 1;} \quad \texttt{x += 1;} \quad \texttt{++x;} \quad \texttt{x++;}

\medskip\textbf{3.3:} a) \texttt{z = x++ + y;} \quad b) \texttt{produto *= 2;} \quad
c) \texttt{produto = produto * 2;} \quad
d) \texttt{if (contador > 10) printf(...);} \quad
e) \texttt{total -= --x;} \quad f) \texttt{total += x--;} \quad
g) \texttt{q \%= divisor;} ou \texttt{q = q \% divisor;} \quad
h) \texttt{123.46} \quad i) \texttt{3.142}

\medskip\textbf{3.6:} \texttt{produto = 25}, \texttt{x = 6}

\medskip\textbf{3.10:} \texttt{z} nunca é decrementado $\Rightarrow$ loop infinito.
Solução: adicionar \texttt{--z;} no corpo do laço.

\end{tcolorbox}

\end{document}
